\documentclass[12pt]{article}

% === Packages ===
\usepackage[margin=1in]{geometry}
\usepackage{amsmath, amssymb}
\usepackage{setspace}
\usepackage{titlesec}
\usepackage{enumitem}
\usepackage{helvet}
\renewcommand{\familydefault}{\sfdefault} % modern sans-serif look

% === Formatting tweaks ===
\setstretch{1.15}
\setlist[enumerate]{itemsep=0.4em, topsep=0.2em}
\titleformat{\section}{\Large\bfseries\sffamily}{\thesection.}{1em}{}
\titleformat{\subsection}{\large\bfseries\sffamily}{\thesubsection}{1em}{}

% === Title ===
\title{\vspace{-2em}\textbf{Applied Mathematics Oral Exam Questions and Model Answers}\\
\large From Euler to Navier--Stokes}
\author{Prepared for Graduate Oral Examination}
\date{}

\begin{document}


\section*{Sample Oral Exam Questionson tensors}

\subsection*{1. What is the material derivative and why is it important in fluid–structure interaction?}

\textbf{Answer.}
The material derivative measures the rate of change of a quantity as experienced
by a moving fluid parcel. For a scalar field $\phi(x,t)$,
\[
\frac{D\phi}{Dt}
= \frac{\partial \phi}{\partial t} + (u\cdot\nabla)\phi,
\]
and for a vector field,
\[
\frac{D\mathbf{u}}{Dt}
= \frac{\partial \mathbf{u}}{\partial t} + (u\cdot\nabla) \mathbf{u}.
\]
This combines local (Eulerian) time dependence with convective effects
(\cite{file:turn4file0}, pp.\ 6–10).

In IBFE, the material derivative appears in:
\begin{itemize}
    \item the \emph{fluid} momentum equation (left-hand side of Navier--Stokes),
    \item the \emph{Lagrangian} update of solid configuration when computing 
          structural velocities and deformation gradients,
    \item coupling the Eulerian fluid velocity to the Lagrangian mesh velocity.
\end{itemize}

It expresses the fact that immersed structures respond to forces determined by
the fluid motion \emph{along their trajectories}.

\subsection*{2. Explain the Cauchy stress tensor and how traction arises from it.}

\textbf{Answer.}
The Cauchy stress tensor $\boldsymbol{\sigma}$ describes how internal surface
forces act within a deforming fluid (\cite{file:turn4file0}, pp.\ 25–32).
For a surface with outward normal $\mathbf{n}$, the traction vector is
\[
\mathbf{t} = \boldsymbol{\sigma}\mathbf{n}.
\]
Each component $\sigma_{ij}$ is the $i$-component of the force acting on a face
whose outward normal is in the $j$-direction.

In IBFE:
\begin{itemize}
    \item This traction represents the \emph{fluid stress} acting on the
          immersed structure.
    \item The traction is used to compute boundary forces on the solid mesh,
          often through surface integrals that feed the structural finite element
          formulation.
    \item Conversely, structural stresses generate elastic forces that are
          spread back to the fluid grid.
\end{itemize}

\subsection*{3. Derive the form of the Navier--Stokes equations used in IBFE.}

\textbf{Answer.}
Starting from Cauchy’s equation of motion (\cite{file:turn4file0}, p.\ 28),
\[
\rho\frac{D\mathbf{u}}{Dt} = \nabla\cdot\boldsymbol{\sigma} + \rho\mathbf{g},
\]
split the stress into pressure and viscous components:
\[
\boldsymbol{\sigma} = -p\mathbf{I} + \boldsymbol{\tau}.
\]
For a Newtonian, incompressible fluid,
\[
\boldsymbol{\tau} = 2\mu \mathbf{D},
\qquad
\mathbf{D} = \tfrac{1}{2}(\nabla \mathbf{u} + \nabla \mathbf{u}^\top)
\]
(\cite{file:turn4file0}, pp.\ 30–33).

This yields
\[
\rho\left( \frac{\partial \mathbf{u}}{\partial t}
+ (\mathbf{u}\cdot\nabla)\mathbf{u} \right)
= -\nabla p + \mu \nabla^2 \mathbf{u} + \rho \mathbf{g},
\qquad 
\nabla\cdot \mathbf{u} = 0.
\]
This is the fluid equation solved in IBFE, with structural forces added as
immersed boundary forcing terms.

\subsection*{4. How does IBFE represent the deformation of an immersed structure?}

\textbf{Answer.}
The solid is described in a Lagrangian coordinate system $X$ with mapping
$\chi(X,t)$ giving its current position.  
The deformation gradient is
\[
\mathbf{F}(X,t) = \nabla_X \chi(X,t).
\]
The solid finite-element formulation provides a constitutive model
(\emph{hyperelastic}, \emph{neo-Hookean}, \emph{fiber-reinforced}, etc.) giving
the first Piola--Kirchhoff stress $\mathbf{P}(X,t)$.

In IBFE:
\begin{itemize}
    \item Lagrangian FE computes internal elastic forces,
    \item these forces are \emph{spread} to the Eulerian fluid grid via IB
          delta kernels,
    \item the fluid velocity is interpolated back to the Lagrangian mesh to
          enforce no-slip (kinematic coupling).
\end{itemize}

\subsection*{5. Explain how the force-spreading step works in IBFE.}

\textbf{Answer.}
Given Lagrangian force density $\mathbf{F}(X,t)$ (often from finite element
quadrature),
the Eulerian force density is
\[
\mathbf{f}(x,t) = \int_{\Omega_s}
\mathbf{F}(X,t)\,\delta(x - \chi(X,t))\,dX.
\]
This maps structural forces to the fluid grid in a conservation-preserving way.
Typically in code, the Dirac delta is replaced by a regularized IB kernel
(e.g., Peskin 4-pt).

This step couples the solid’s internal stresses to the fluid’s momentum
equation.

\subsection*{6. What does the no-slip condition mean in IBFE?}

\textbf{Answer.}
IBFE enforces
\[
\frac{\partial \chi}{\partial t}(X,t) = \mathbf{u}(\chi(X,t),t),
\]
i.e.,
the solid moves with the local fluid velocity.
This is implemented by \emph{interpolation} from Eulerian velocity to the
Lagrangian mesh via the same IB kernel used for force spreading.

No explicit boundary conditions are applied on the Eulerian grid;
instead, structural forces modify the fluid so that this kinematic condition
is satisfied.

\subsection*{7. What conditions ensure that a matrix of components represents a tensor?}

\textbf{Answer.}
A $3\times 3$ array $A_{ij}$ represents a tensor if it obeys the tensor
transformation law under a change of orthonormal basis:
\[
A^{(m)} = Q A^{(e)} Q^T.
\]
If it does not transform this way, it is just a matrix.
This is essential in IBFE because deformation gradients, stresses, and rate-of-strain
tensors must be genuine geometric objects independent of the coordinate
representation.

\subsection*{8. Why must the viscous stress tensor be symmetric?}

\textbf{Answer.}
In fluids without couple stresses,
\[
T_{ij} = T_{ji}.
\]
Symmetry ensures angular momentum conservation at the infinitesimal level.
This is needed for the derivation of Navier--Stokes from Cauchy momentum
equation.

In IBFE, symmetric stress ensures that structural forces do not artificially
inject or remove angular momentum from the fluid.

\end{document}